\documentclass[10pt,a4paper]{amsart}
\usepackage[margin=19mm]{geometry}
\usepackage{amsmath,amssymb,mathtools}
\usepackage[scr=rsfs,cal=euler]{mathalfa}
\usepackage{xfrac}
\usepackage[unicode,hidelinks,bookmarks=false]{hyperref}
\newcommand{\ie}{i.e.\ }
\newcommand{\cf}{cf.\ }
\newcommand{\resp}{respectively\ }
\newcommand{\Subsection}{\subsection}
\providecommand{\nobreakdash}{\nobreak\hbox{-}}
\setlength{\emergencystretch}{2em}
\multlinegap0pt
\usepackage{bm}
\usepackage{bbm}
\usepackage{dsfont}
\usepackage{xspace}
\usepackage{cancel}
\usepackage{mathtools}
\usepackage{amsthm}
\makeatletter
\@ifundefined{theorem}{\newtheorem{theorem}{Theorem}}{}
\makeatother
\title[Control semiflows, chain controllability, and the Selgrade d]{Control semiflows, chain controllability, and the Selgrade decomposition for linear delay systems}
\author{Fritz Colonius}
\date{Source: arXiv:2408.10048v1 (CC BY 4.0)}
\begin{document}
\maketitle

\noindent\emph{Source and licence.} Control semiflows, chain controllability, and the Selgrade decomposition for linear delay systems. Version arXiv:2408.10048v1 (CC BY 4.0), distributed under the Creative Commons Attribution 4.0 International licence, as recorded in the arXiv licence metadata for this exact version and shown on the abstract page https://arxiv.org/abs/2408.10048v1. Full rights notice: licenses/arxiv-2408.10048-CC-BY-4.0.txt.\par\smallskip

\noindent\textit{Editorial scope.} Counted unit: the Theorem whose source label is \texttt{Theorem\_ccs1} in the article (source0.tex lines 1165--1169, source label \texttt{Theorem\_ccs1}), delivered together with the article's own complete original proof of it (source0.tex lines 1171--1374, one \texttt{proof} environment of 204 lines). No mathematics is written, repaired or re-derived: statement and proof are the article's own text line for line. Required context retained uncounted: D1 (lines 67--72); mild (lines 511--514). Every definition, display and label of those lines is retained unchanged. Omitted from this package and not redistributed: 1--66, 73--510, 515--1164, 1170--1170, 1375--2330. Nothing inside a retained excerpt is omitted or altered: every retained body is delivered exactly as the article's lines are written, so no omission receipt of any kind accompanies this package. Citations: the retained excerpts carry no citation command at all, so no bibliography entry of the article is transcribed and no reference is redistributed. Reference adaptation: none was needed and none was made; every reference inside the delivered unit resolves inside it, so no unresolved reference or citation is produced. Printed slips kept literal: no printed word, symbol, hypothesis or number of the counted statement or of its proof was altered. Layout and class: the article's own class is \texttt{siamltex} with packages including amsfonts, amsmath, amssymb, graphicx, enumerate, color, hyperref; the wrapper is the lane's pinned \texttt{amsart} editorial wrapper at 10pt on A4, which restates the theorem-like environments the retained text needs and adds the article's own macro definitions that the retained text needs (none), with extra wrapper packages (bm, bbm, dsfont, xspace, cancel, mathtools). The delivered numbering is the unit's own. Assets: no retained span contains a figure environment, a graphics-inclusion command, a PDF-page inclusion, a table or a TikZ picture; no image file is copied into this package, the package declares no asset dependency and no image file is redistributed with it. Licence: version arXiv:2408.10048v1 of this article is distributed under the Creative Commons Attribution 4.0 International licence, as recorded in the arXiv licence metadata for this exact version and shown on the abstract page https://arxiv.org/abs/2408.10048v1. A full rights notice is delivered at licenses/arxiv-2408.10048-CC-BY-4.0.txt. Characters: every retained excerpt is pure ASCII, so the delivered engine, XeLaTeX with its default Latin Modern text font, reports no missing character.
\begin{align}
\dot{x}(t)  &  =A_{0}x(t)+\sum_{i=1}^{p}A_{i}x(t-h_{i})+B_{0}u(t)+\sum
_{i=1}^{p}B_{i}u(t-h_{i}),\,u\in\mathcal{U},\label{D1}\\
x(0)  &  =r,\,x(s)=f(s)\text{ for almost all }-h\leq s\leq0\text{ and
}u(t)=0\text{ for }t\leq0.\nonumber
\end{align}

\begin{equation}
y(t)=T(t)y_{0}+\int_{0}^{t}T(t-s)B\left(  u(s),u(s-h_{1}),\ldots
,u(s-h_{p})\right)  ds. \label{mild}%
\end{equation}

\begin{theorem}
\label{Theorem_ccs1}Consider the linear delay control system given by
(\ref{D1}) and assume that the matrix $A_{p}$ is invertible. Then there exists
a unique chain control set $E$ in $M_{2}$.
\end{theorem}

\begin{proof}
First note that for $u\equiv0$ the origin $0\in M_{2}$ is an equilibrium,
hence $\{0\}$ is a weakly invariant chain controllable set. Define $E$ as the
union of all weakly invariant chain controllable sets containing $\{0\}$. Then
$E$ is a weakly invariant chain controllable set and certainly it is maximal
with these properties, hence it is a chain control set. It remains to prove uniqueness.

Observe that the trajectories $y(t)=\varphi(t,y_{0},u),t\in\mathbb{R}$, of
(\ref{mild}) satisfy, for $\alpha\in(0,1)$,%
\begin{align}
\alpha\varphi(t,y_{0},u)  &  =\alpha T(t)y_{0}+\alpha\int_{0}^{t}%
T(t-s)B\left(  u(s),u(s-h_{1}),\ldots,u(s-h_{p})\right)  ds\nonumber\\
&  =T(t)\alpha y_{0}+\int_{0}^{t}T(t-s)B\left(  \alpha u(s),\alpha
u(s-h_{1}),\ldots,\alpha u(s-h_{p})\right)  ds\nonumber\\
&  =\varphi(t,\alpha y_{0},\alpha u). \label{alpha1}%
\end{align}
Here, $\varphi(\cdot,\alpha y_{0},\alpha u)$ is a trajectory of (\ref{mild})
since $\Omega$ is a convex neighborhood of $0\in\mathbb{R}^{m}$ implying that
the controls $\alpha u$ are in $\mathcal{U}$.

Suppose that $E^{\prime}$ is any chain control set and let $y\in E^{\prime}$.
First we will construct controlled $(\varepsilon,\tau)$-chains from $y$ to
$0\in E$.

\textbf{Step 1:} There is a controlled $(\varepsilon,\tau)$-chain from $y$ to
$\alpha y$ for some $\alpha\in(0,1)$.

For the proof consider a controlled $(\varepsilon/2,\tau)$-chain $\zeta$ from
$y$ to $y$ given by $y_{0}=y,y_{1},\ldots,y_{q}=y,\,u_{0},\ldots,u_{q-1}%
\in\mathcal{U}$, and $\tau_{0},\ldots,\tau_{q-1}\geq\tau$ with%
\[
\left\Vert \varphi(\tau_{i},y_{i},u_{i})-y_{i+1}\right\Vert <\varepsilon
/2\text{ for }i=0,\ldots,q-1.
\]
Let $\alpha\in(0,1)$ with $(1-\alpha)\left\Vert y\right\Vert <\varepsilon/2$.
It follows that%
\[
\left\Vert \varphi(\tau_{q-1},y_{q-1},u_{q-1})-\alpha y_{q}\right\Vert
\leq\left\Vert \varphi(\tau_{q-1},y_{q-1},u_{q-1})-y\right\Vert +\left\Vert
y-\alpha y\right\Vert <\varepsilon.
\]
This defines a controlled $(\varepsilon,\tau)$-chain $\zeta^{(1)}$ from $y$ to
$\alpha y$.

\textbf{Step 2: }Replacing $y_{i}$ by $\alpha y_{i}$ and $u_{i}$ by $\alpha
u_{i}$ we get by (\ref{alpha1})%
\begin{align*}
\left\Vert \varphi(\tau_{i},\alpha y_{i},\alpha u_{i})-\alpha y_{i+1}%
\right\Vert  &  =\alpha\left\Vert \varphi(\tau_{i},y_{i},u_{i})-y_{i+1}%
\right\Vert <\varepsilon/2\text{ for }i=0,\ldots,q-1,\\
\left\Vert \varphi(\tau_{q-1},\alpha y_{q-1},\alpha u_{q-1})-\alpha
^{2}y\right\Vert  &  \leq\left\Vert \varphi(\tau_{q-1},\alpha y_{q-1},\alpha
u_{q-1})-\alpha y\right\Vert +\left\Vert \alpha y-\alpha^{2}y\right\Vert
<\varepsilon.
\end{align*}
This defines a controlled $(\varepsilon,\tau)$-chain $\zeta^{(2)}$ from
$\alpha y$ to $\alpha^{2}y$. The concatenation of $\zeta^{(2)}$ and
$\zeta^{(1)}$ yields a controlled $(\varepsilon,\tau)$-chain $\zeta^{(2)}%
\circ\zeta^{(1)}$ from $y$ to $\alpha^{2}y$.

Repeating this construction, we find that the concatenation $\zeta^{(k)}%
\circ\cdots\circ\zeta^{(1)}$ is a controlled $(\varepsilon,\tau)$-chain from
$y\in E^{\prime}$ to $\alpha^{k}y$. Since $\alpha^{k}\rightarrow0$ for
$k\rightarrow\infty$, we can take $k\in\mathbb{N}$ large enough, such that the
last piece of the chain $\zeta^{(k)}$ satisfies%
\[
\left\Vert \varphi(\tau_{q-1},a^{k}y_{q-1},a^{k}u_{q-1})\right\Vert
<\varepsilon.
\]
Thus we may take $0\in E$ as the final point of this controlled $(\varepsilon
,\tau)$-chain showing that the concatenation $\zeta^{(k)}\circ\cdots\circ
\zeta^{(1)}$ defines a controlled $(\varepsilon,\tau)$-chain from $y\in
E^{\prime}$ to $0\in E$.

\textbf{Step 3:} Next we construct controlled chains from $0$ to $y\in
E^{\prime}$.

Consider a controlled $(\varepsilon,\tau)$-chain from $y$ to $y$ given by
$y_{0}=y,y_{1},\ldots,y_{q}=y,\allowbreak\,u_{0},\ldots,u_{q-1}\in\mathcal{U}%
$, and $\tau_{0},\ldots,\tau_{q-1}\geq\tau$ with%
\[
\left\Vert \varphi(\tau_{i},y_{i},u_{i})-y_{i+1}\right\Vert <\varepsilon\text{
for }i=0,\ldots,q-1.
\]
For every $\alpha\in(0,1)$ formula (\ref{alpha1}) shows that $\alpha
y_{0}=\alpha y,\alpha y_{1},\ldots,\alpha y_{q}=\alpha y,\,\allowbreak\alpha
u_{0},\ldots,\allowbreak\alpha u_{q-1}\in\mathcal{U}$, and $\tau_{0}%
,\ldots,\tau_{q-1}\geq\tau$ define a controlled $(\alpha\varepsilon,\tau)$-
chain from $\alpha y$ to $\alpha y$ with%
\[
\left\Vert \varphi(\tau_{i},\alpha y_{i},\alpha u_{i})-\alpha y_{i+1}%
\right\Vert <\alpha\varepsilon\text{ for }i=0,\ldots,q-1.
\]
Let $\alpha\in(0,\varepsilon)$ be small enough such that $\alpha\left\Vert
y\right\Vert <\varepsilon$. Then we may add a segment $\varphi(t,0,0)=0,t\in
\lbrack0,\tau]$, at the beginning to obtain a controlled $(\alpha
\varepsilon,\tau)$-chain from $0$ to $\alpha y$. Furthermore, we find%
\begin{align*}
\left\Vert \varphi(\tau_{q-1},\alpha y_{q-1},\alpha u_{q-1})-(\alpha
+\varepsilon)y\right\Vert  &  \leq\left\Vert \varphi(\tau_{q-1},\alpha
y_{q-1},\alpha u_{q-1})-\alpha y\right\Vert +\varepsilon\left\Vert
y\right\Vert \\
&  \leq\alpha\varepsilon+\varepsilon\left\Vert y\right\Vert <(\varepsilon
+\left\Vert y\right\Vert )\varepsilon.
\end{align*}
Taking $\varepsilon\in(0,1)$ we have constructed a controlled $((1+\varepsilon
)\left\Vert y\right\Vert ,\tau)$-chain $\zeta^{(1)}$ from $0$ to
$(\alpha+\varepsilon)y$.

Now construct a controlled $(2\varepsilon,\tau)$-chain $\zeta^{(2)}$ from
$(\alpha+\varepsilon)y$ to $(2\alpha+\varepsilon)y$. By (\ref{alpha1}),%
\[
\left\Vert \varphi(\tau_{i},(\alpha+\varepsilon)y_{i},(\alpha+\varepsilon
)u_{i})-(\alpha+\varepsilon)y_{i+1}\right\Vert <(\alpha+\varepsilon
)\varepsilon\text{ for }i=0,\ldots,q-1,
\]
and, since $\alpha\left\Vert y\right\Vert <\varepsilon$,%
\begin{align*}
&  \left\Vert \varphi(\tau_{q-1},(\alpha+\varepsilon)y_{q-1},(\alpha
+\varepsilon)u_{q-1})-(2\alpha+\varepsilon)y\right\Vert \\
&  \leq\left\Vert \varphi(\tau_{q-1},(\alpha+\varepsilon)y_{q-1}%
,(\alpha+\varepsilon)u_{q-1})-(\alpha+\varepsilon)y\right\Vert +\alpha
\left\Vert y\right\Vert \\
&  <(\alpha+\varepsilon)\varepsilon+\varepsilon.
\end{align*}
For $\alpha+\varepsilon<1$, it follows that $(\alpha+\varepsilon
)\varepsilon+\varepsilon\leq2\varepsilon$ and hence this is a controlled
$(2\varepsilon,\tau)$-chain from $(\alpha+\varepsilon)y$ to $(2\alpha
+\varepsilon)y$.

Next construct a controlled $(2\varepsilon,\tau)$-chain $\zeta^{(3)}$ starting
in $(2\alpha+\varepsilon)y$: By (\ref{alpha1}),%
\[
\left\Vert \varphi(\tau_{i},(2\alpha+\varepsilon)y_{i},(2\alpha+\varepsilon
)u_{i})-(2\alpha+\varepsilon)y_{i+1}\right\Vert <(2\alpha+\varepsilon
)\varepsilon\text{ for }i=0,\ldots,q-1
\]
and%
\begin{align*}
&  \left\Vert \varphi(\tau_{q-1},(2\alpha+\varepsilon)y_{q-1},(2\alpha
+\varepsilon)u_{q-1})-(3\alpha+\varepsilon)y\right\Vert \\
&  \leq\left\Vert \varphi(\tau_{q-1},(2\alpha+\varepsilon)y_{q-1}%
,(2\alpha+\varepsilon)u_{q-1})-(2\alpha+\varepsilon)y\right\Vert
+\alpha\left\Vert y\right\Vert \\
&  <(2\alpha+\varepsilon)\varepsilon+\varepsilon.
\end{align*}
For $2\alpha+\varepsilon<1$, this defines a controlled $(2\varepsilon,\tau
)$-chain from $(2\alpha+\varepsilon)y$ to $(3\alpha+\varepsilon)y$.

As long as $j\alpha+\varepsilon<1$, we can proceed in this way to obtain
controlled $(2\varepsilon,\tau)$-chains $\zeta^{(j+1)}$ from $(j\alpha
+\varepsilon)y$ to $((j+1)\alpha+\varepsilon)y$ satisfying%
\[
\left\Vert \varphi(\tau_{i},(j\alpha+\varepsilon)y_{i},(j\alpha+\varepsilon
)u_{i})-(j\alpha+\varepsilon)y_{i+1}\right\Vert <(j\alpha+\varepsilon
)\varepsilon<\varepsilon\text{ for }i=0,\ldots,q-1,
\]
with%
\begin{align*}
&  \left\Vert \varphi(\tau_{q-1},(j\alpha+\varepsilon)y_{q-1},(j\alpha
+\varepsilon)u_{q-1})-((j+1)\alpha+\varepsilon)y\right\Vert \\
&  \leq\left\Vert \varphi(\tau_{q-1},(j\alpha+\varepsilon)y_{q-1}%
,(j\alpha+\varepsilon)u_{q-1})-(j\alpha+\varepsilon)y\right\Vert
+\alpha\left\Vert y\right\Vert \\
&  \leq(j\alpha+\varepsilon)\varepsilon+\alpha\left\Vert y\right\Vert
<2\varepsilon.
\end{align*}
When for $j=k$ we arrive at $k\alpha+\varepsilon<1$ and $(k+1)\alpha
+\varepsilon\geq1$, we find%
\begin{equation}
\varepsilon>\varepsilon+k\alpha+\varepsilon-1>(k+1)\alpha+\varepsilon-1\geq0.
\label{N+1}%
\end{equation}
Thus we get for $i=0,\ldots,q-1$%
\[
\left\Vert \varphi(\tau_{i},(k\alpha+\varepsilon)y_{i},(k\alpha+\varepsilon
)u_{i})-(k\alpha+\varepsilon)y_{i+1}\right\Vert <(k\alpha+\varepsilon
)\varepsilon<\varepsilon,
\]
and, by (\ref{N+1}),%
\begin{align*}
&  \left\Vert \varphi(\tau_{q-1},(k\alpha+\varepsilon)y_{q-1},(k\alpha
+\varepsilon)u_{q-1})-y\right\Vert \\
&  \leq\left\Vert \varphi(\tau_{q-1},(k\alpha+\varepsilon)y_{q-1}%
,(k\alpha+\varepsilon)u_{q-1})-((k+1)\alpha+\varepsilon)y\right\Vert \\
&  \qquad+\left\Vert ((k+1)\alpha+\varepsilon)y-y\right\Vert \\
&  <\left\Vert \varphi(\tau_{q-1},(k\alpha+\varepsilon)y_{q-1},(k\alpha
+\varepsilon)u_{q-1})-(k\alpha+\varepsilon)y\right\Vert +\varepsilon\left\Vert
y\right\Vert \\
&  \qquad+\left\Vert ((k+1)\alpha+\varepsilon-1)y\right\Vert \\
&  <(k\alpha+\varepsilon)\varepsilon+\varepsilon\left\Vert y\right\Vert
+((k+1)\alpha+\varepsilon-1)\left\Vert y\right\Vert \\
&  <\varepsilon+\varepsilon\left\Vert y\right\Vert +\varepsilon\left\Vert
y\right\Vert <\varepsilon+2\varepsilon\left\Vert y\right\Vert .
\end{align*}
Thus this defines a controlled $((1+2\left\Vert y\right\Vert )\varepsilon
,\tau)$-chain $\zeta^{(k+1)}$ from $(k\alpha+\varepsilon)y$ to $y$. The
concatenation $\zeta^{(k+1)}\circ\zeta^{(k)}\circ\cdots\circ\zeta^{(1)}$
yields a controlled $((1+2\left\Vert y\right\Vert )\varepsilon,\tau)$-chain
from $0$ to $y$.

Since $\varepsilon,\tau>0$ are arbitrary, steps 2 and 3 imply that $y\in
E^{\prime}\cap E$ and hence $E^{\prime}=E$.
\end{proof}

\end{document}
